\documentclass[11pt,oneside]{scrbook}

% --- Packages ---
\usepackage[utf8]{inputenc}
\usepackage[T1]{fontenc}
\usepackage{geometry}
\geometry{a4paper, margin=1in, headheight=25pt}
% \usepackage{newtxtext,newtxmath} % Removed custom Times New Roman font to use default Computer Modern
\usepackage{xcolor}
\usepackage{graphicx}
\usepackage{booktabs}
\usepackage{array}
\usepackage{hyperref}
\usepackage{scrlayer-scrpage} % For headers and footers
\usepackage{amsmath}
\usepackage{caption}
\usepackage{longtable}
\usepackage{enumitem}

% --- Header and Footer Configuration ---
\clearpairofpagestyles
\ihead{\headmark} % Chapter title on the left
\ohead{\pagemark}  % Page number on the right
\automark{chapter}
\setkomafont{pagehead}{\normalfont\itshape\small}
\setkomafont{pagenumber}{\normalfont\bfseries}

% Adding a horizontal rule for the "top-bar" feel
\KOMAoptions{headsepline=0.5pt} % Modern way to set header line thickness

% --- Formatting ---
\definecolor{navy}{RGB}{0, 31, 63}
\hypersetup{colorlinks=true, linkcolor=navy, urlcolor=blue, citecolor=navy}
\setlist[description]{style=nextline}

% --- Custom Macros ---
\newcommand{\projecttitle}{Agentic Sentiment Intelligence System}
\newcommand{\mysubtitle}{Bridging the Sarcasm Gap in E-commerce via Multi-Tiered Neural Architectures}

% --- Title Page ---
\title{
    \vspace{2cm}
    \Huge \textbf{\projecttitle}\\
    \Large \textit{\mysubtitle}\\
    \vspace{1.5cm}
    \Large Capstone Project Final Report (2026)
}
\author{
    Lokesh Patra | 020 \\
    Aditya Narayan Purohit | 006 \\
    Aakash Mahato | 016 \\
    Rudra Narayan Behera | 017 \\
    Subham Nayak | 019
}
\date{April 2026}

\begin{document}

\frontmatter
\maketitle

\chapter{Executive Summary}
The \textbf{Agentic Sentiment Intelligence System} is a state-of-the-art solution designed to bridge the "Sarcasm Gap" in e-commerce review analysis. Traditional sentiment models---while effective at identifying polarized language---frequently fail when faced with irony, nuanced negativity, and phonetic slang. This project introduces a \textbf{Multi-Tiered Agentic Architecture} that combines neural normalization, transformer-based classification, and symbolic reasoning. By integrating an \textbf{Active Data Mining} workflow, the system identifies "Ironic Dissonance" (conflicts between star ratings and text) to train a specialist engine capable of high-precision sentiment detection in complex scenarios.

\chapter{Abstract}
This project presents an \textbf{Agentic Sentiment Intelligence System} designed to solve the critical "Sarcasm Gap" in e-commerce review analysis. Traditional models are adept at polarized language, they frequently misclassify reviews featuring "Ironic Dissonance" (e.g., \textit{"This phone is so great at being terrible"}). Our system introduces a Multi-Tiered Agentic Architecture combining neural normalization via a \textbf{Linguistic Quality Node (LQN)}, transformer-based classification using a fine-tuned \textbf{DistilRoBERTa Specialist}, and an \textbf{Agentic Reasoning Engine (RE)} for symbolic explanation. By integrating an Active Data Mining workflow to identify and weight ironic samples, the system achieves over 87\% accuracy on complex, sarcasm-heavy datasets, all while optimized for CPU-only hardware democratization.

\tableofcontents
\listoftables
\listoffigures

\mainmatter

\chapter{Project Vision \& Problem Statement}

\section{The Strategic Problem: The Sarcasm Gap}
The "Sarcasm Gap" represents a critical failure point in modern e-commerce sentiment analysis. While traditional models identify obvious polarized words, they frequently falter when encountering "Ironic Dissonance"---a phenomenon where the semantic content of a review contradicts its intended sentiment. 

Standard "Bag-of-Words" models (Logistic Regression, Random Forest) are fundamentally limited by their reliance on keyword frequency and ignore bidirectional relationship and negative modifiers (e.g., weighting "perfectly" as positive even when followed by "mis-crafted"). When sarcasm is misclassified, business intelligence is compromised, leading to skewed data that masks severe product defects.

\section{Core Objectives}
\begin{enumerate}
    \item \textbf{Solve the Sarcasm Gap:} Move beyond keyword matching to understanding bidirectional sequence context.
    \item \textbf{Neural Normalization:} Automatically "clean" and "rewrite" messy, slang-heavy user reviews before classification using generative models (LQN).
    \item \textbf{Agentic Reasoning:} Provide human-readable explanations ("Why") for every prediction to improve transparency and trust.
    \item \textbf{Hardware Democratization:} Ensure the system runs efficiently on consumer-grade hardware (8GB RAM / CPU-only) without requiring GPUs.
\end{enumerate}

\chapter{Literature Survey \& Gap Analysis}

\section{Comprehensive Literature Survey}
The following table provides a systematic review of the literature from 2021 to 2026, mapping contemporary research against the Agentic Sentiment Intelligence System.

\begin{longtable}{p{2.2cm}p{2.5cm}p{2.5cm}p{3.2cm}p{3.2cm}}
\caption{Comprehensive Literature Survey} \\
\toprule
\textbf{Author \& Year} & \textbf{Title} & \textbf{Contributions} & \textbf{Literature Gap} & \textbf{Project Fix} \\
\midrule
\endfirsthead
\toprule
\textbf{Author \& Year} & \textbf{Title} & \textbf{Contributions} & \textbf{Literature Gap} & \textbf{Project Fix} \\
\midrule
\endhead
\bottomrule
\endfoot
Maragheh et al. (2025) & ARAG: Agentic RAG & Multi-agent collaboration in RAG pipelines. & Fails to capture nuanced preferences. & Multi-Tiered Agentic Architecture. \\
\addlinespace
Cao \& Hu (2026) & Solicit-Then-Suggest Agentic Purchasing & Multi-round conversations for preference. & Lacks ability to learn from messy descriptions. & LQN to programmatically repair descriptions. \\
\addlinespace
Ayanami0730 (2026) & A-RAG: Scaling Agentic RAG & Hierarchical tools for agent autonomy. & Rigid RAG systems without model participation. & Intelligence pipeline (Normalize $\to$ Analyze $\to$ Explain). \\
\addlinespace
Scott et al. (2025) & ABSA on Amazon Reviews & Novel Hybrid ML Algorithm (SVHA). & struggles with complex language structure. & Sublinear TF-IDF and neural normalization. \\
\addlinespace
Wen \& Rezapour (2026) & Transformer \& Prototype Model & Prototype tensors for sentence-level XAI. & Word-level XAI is non-intuitive for irony. & RE providing symbolic "Conflict Logic". \\
\addlinespace
Daraghmi \& Zyadeh (2026) & Sentiment Classification study & Benchmarks ML and DL models. & Fails to capture sarcasm and negation. & Contextual Persistence in DistilRoBERTa. \\
\addlinespace
Thogesan et al. (2025) & XAI with LLMs & Layer-wise SHAP analysis for BERT. & LLMs function as "black boxes". & RE mapping confidence to insights. \\
\addlinespace
Dubey et al. (2025) & Unpacking Sarcasm & Contextual Transformer Approach. & Lack of structured preprocessing. & Latency reduction 16s to 1.1s via pruning. \\
\addlinespace
Akter et al. (2025) & Sentiment Analysis Impact & Validates BERT outperforms ML. & High computational resource demands. & Hardware Democratization (8GB RAM/CPU). \\
\addlinespace
Agrawal (2025) & Future of E-Commerce AI & Specialized agents for life domains. & Lacks goal-oriented understanding. & End-to-end reasoning pipeline. \\
\addlinespace
Trivedi (2025) & Explainable Agentic AI & Conceptual models for agentic search. & Transparency issues in autonomy. & RE providing verifiable reasoning trace. \\
\addlinespace
Chauhan \& Kumar (2024) & Multilingual Sarcasm & mBERT for code-mixed text. & Emojis and slang degrade accuracy. & Emoji Demojization and custom slang map. \\
\addlinespace
Kaseb \& Farouk (2023) & SAIDS: Dialect and Sarcasm & Informed decision architecture. & Literal vs Intended meaning conflict. & Active Mining of Ironic Dissonance. \\
\addlinespace
Abu Farha et al. (2022) & Sarcasm Detection easy? & Benchmarks human vs machine. & Machine models struggle with context. & Specialist Engine trained on "Hard Cases". \\
\addlinespace
Romio, G. (2026) & 15 Best AI SA Tools & Reviews top-tier platforms for 2026. & Basic keyword matching false alarms. & RE flags "Conflict of Signals". \\
\end{longtable}

\chapter{Research Methodology \& Timeline}

\section{The Modified CRISP-DM Approach}
We followed a 6-phase scientific approach to ensure data integrity:
\begin{description}
    \item[Phase 1] Problem Definition \& Data Selection (Amazon, Flipkart, PUMA).
    \item[Phase 2] Data Refinery \& Feature Engineering (Regex cleaning, Lemmatization).
    \item[Phase 3] Baseline Modeling (LR, RF, XGBoost).
    \item[Phase 4] Synthesis \& Deep Intelligence (BERT/RoBERTa).
    \item[Phase 5] Implementation \& Visual Analytics (FastAPI, Streamlit).
    \item[Phase 6] Final Evaluation \& Documentation.
\end{description}

\section{The 10-Week Evolution}
The project followed a rigorous 10-week development cycle, evolving from a standard statistical classifier into a sophisticated neural agent.
\begin{description}
    \item[Weeks 1--2: Initiation \& Diagnostic] Analyzed 200k+ rows; identified heavy imbalance. Resolved scikit-learn and XGBoost mismatches that caused crashes during high-dimensional matrix operations.
    \item[Weeks 3--4: Data Refinery \& Classical Modeling] Trained RF and XGBoost. Reached 84\% accuracy but confirmed systematic sarcasm failure. Applied SMOTE to address Neutral class deficiency.
    \item[Weeks 5--6: The Context Revolution] Pivoted to Transformers. Discovered "Neural Echoing" in T5 models where models repeated input instead of repairing it.
    \item[Weeks 7--8: Agentic System Development] Architected the LQN and RE; designed the professional, emoji-free Glassmorphism Dashboard.
    \item[Weeks 9--10: Specialist Optimization \& Final Polish] Engineered the Ironic Dissonance Miner; fine-tuned DistilRoBERTa Specialist; reduced latency from 16.7s to 1.1s.
\end{description}

\chapter{Data Refinery \& Preprocessing}

\section{Intricate Preprocessing Pipeline}
The system employs a "Deep Clean" strategy to preserve sentimental signals:
\begin{itemize}
    \item \textbf{Fuzzy Column Mapping:} Implemented logic that programmatically detects headers ('Summary', 'Review'), ensuring the pipeline remains robust regardless of data source.
    \item \textbf{Emoji Demojization:} Converts icons into readable text (\texttt{:smiling\_face:}) to capture emotive tone.
    \item \textbf{Phonetic Slang Mapping:} Maps shorthand (e.g., 'gr8' $\to$ 'great') to reduce vocabulary for improved model focus.
    \item \textbf{TF-IDF Tuning:} Applied \texttt{sublinear\_tf=True} to dampen high-frequency but low-value words.
\end{itemize}

\section{Dataset Imbalance \& SMOTE}
The initial 200k+ dataset (Dataset-SA.csv) was heavily biased. We applied \textbf{SMOTE (Synthetic Minority Over-sampling Technique)} to balance classes, which improved XGBoost's \textbf{Neutral Recall to 58.93\%}.

\chapter{System Architecture \& Agentic Design}

\section{The Agentic Intelligence Pipeline}
Our system treats sentiment as a three-stage intelligence pipeline rather than a single classification task:
\begin{enumerate}
    \item \textbf{Normalize (LQN):} The language is synthesized into standard English.
    \item \textbf{Analyze (Specialist Engine):} The context is evaluated for irony.
    \item \textbf{Explain (Reasoning Engine):} The rationale is provided to the user.
\end{enumerate}

\section{Three-Phase Agentic Development Workflow}
A decoupled, tiered architecture is essential for scaling AI services across heterogeneous environments. By separating the layers, the system follows a 3-phase progression:
\begin{enumerate}
    \item \textbf{Empirical Baseline:} Establishing a performance ceiling using lightweight classical models (LR, RF, XGB) to understand basic dataset patterns.
    \item \textbf{Neural Integration:} Transitioning to Deep Learning via BERT-base to capture natural language context and identifying "Intensity Bias" (being fooled by positive keywords).
    \item \textbf{Specialist Active Mining:} Utilizing "Active Learning" to identify ironic dissonance and training a specialist DistilRoBERTa model to resolve it.
\end{enumerate}

\section{The Agentic Architecture: LQN + RE}
The uniqueness of this project lies in these "Intelligent Buffers":
\begin{description}
    \item[A. Linguistic Quality Node (LQN)] A specialized Transformer (\texttt{pszemraj/grammar-synthesis-base}) fine-tuned for grammar synthesis. It rewrites input before it reaches the classifier. Optimization includes \textbf{Beam Search (k=5)} and \textbf{Early Stopping} to ensure repairs occur in under 400ms on a standard CPU.
    \item[B. Reasoning Engine (RE)] Provides Explainable AI (XAI) by mapping neural confidence scores and linguistic signals into actionable reports.
\end{description}

\section{System Design Functional Map}
\begin{longtable}{p{2.5cm}p{3.5cm}p{7.5cm}}
\toprule
\textbf{Layer} & \textbf{Component} & \textbf{Strategic Value} \\
\midrule
Data Mining & Active Sarcasm Miner & Identifies "Hard Cases" via dissonance filtering. \\
LQN & Neural Normalizer & Denoises phonetic slang and character-level errors. \\
Specialist & DistilRoBERTa & High-precision context-aware sentiment detection. \\
Backend & FastAPI Service & Asynchronous pipeline serving for concurrent requests. \\
Frontend & Glassmorphism UI & Professional dashboard with Explainable AI (XAI). \\
\bottomrule
\end{longtable}

\chapter{The Golden Rule: Active Sarcasm Mining}

\section{Specialist Active Mining (Ironic Dissonance)}
To solve sarcasm without manual labeling, we developed a breakthrough workflow:
\begin{itemize}
    \item \textbf{The Detector:} Scanned 205,052 reviews for "Dissonance"---where Star Rating was 1-2 but text contained Positive Intensifiers ("Wow", "Amazing").
    \item \textbf{Correction:} Programmatically identified 336 "Pure Sarcasm" cases and forced them to 'Negative'.
    \item \textbf{Synthetic Augmentation:} Applied \textbf{10x Weighting} during training to ensure patterns were learned.
    \item \textbf{Result:} Final model (DistilRoBERTa) learned that "Wow" plus a low rating signifies extreme disappointment.
\end{itemize}

\chapter{Experimental Results \& Ablation Studies}

\section{Classical ML Modeling \& Evaluation}
We trained three primary baseline models to establish an empirical ceiling:
\begin{itemize}
    \item \textbf{Logistic Regression (LR):} Reached 81.8\% accuracy; high-speed sanity check.
    \item \textbf{Random Forest (RF):} Overall champion for simple reviews (84.6\%). However, model file swelled to 1.5 GB, leading to \textbf{hardware bottlenecks} during deployment.
    \item \textbf{XGBoost (XGB):} Neutral specialist; 58.9\% recall on notoriously difficult neutral reviews.
\end{itemize}

\section{Model Comparisons}
\begin{table}[h]
\centering
\begin{tabular}{p{4.5cm}p{2.5cm}p{8cm}}
\toprule
\textbf{Model Attempted} & \textbf{Accuracy} & \textbf{Strategic Strength} \\
\midrule
Random Forest & 84.6\% & High speed; optimal for simple linear reviews. \\
XGBoost & 80.9\% & Neutral specialist; high recall on 'hard-to-classify' rows. \\
BERT-Base (General) & 82.1\% & Context-aware but lacked specific irony training. \\
\textbf{DistilRoBERTa Specialist} & \textbf{87.0\%} & \textbf{Contextual Persistence}; deep irony understanding. \\
\bottomrule
\end{tabular}
\caption{Inference Model Performance Comparison}
\end{table}

\section{LQN Ablation Study (Pass and Failures)}
\begin{table}[h]
\centering
\begin{tabular}{p{4cm}p{3cm}p{8cm}}
\toprule
\textbf{Neural Model} & \textbf{Result} & \textbf{Technical Rationale} \\
\midrule
FLAN-T5-Base & Rejected & "Neural Echoing" (repeated input without repairs). \\
BART-Large-CNN & Rejected & News Hallucinations and 6.2GB RAM load. \\
\textbf{Grammar-Synthesis} & \textbf{Success} & Final Production; optimal for CPU latency (1.1s). \\
\bottomrule
\end{tabular}
\caption{Neural Normalization Ablation Study}
\end{table}

\chapter{Implementation \& Optimization}

\section{FastAPI Infrastructure Layer}
The backend is powered by a FastAPI service layer, chosen for its asynchronous, decoupled design which is ideal for real-time inference.
\begin{itemize}
    \item \textbf{Request/Response Format:}
    \begin{itemize}
        \item Request: \texttt{\{"text": "...", "model\_choice": "..."\}}
        \item Response: \texttt{\{"sentiment": "...", "confidence": 0.89, "clean\_text": "..."\}}
    \end{itemize}
    \item \textbf{Initialization:} The system's \textbf{12-second startup time} is primarily attributed to the loading of the 1.5 GB Random Forest model weights.
\end{itemize}

\section{Frontend Design (Glassmorphism)}
The dashboard adopts a "Premium Professional Aesthetic" following a strict **Emoji-Free Protocol**. 
\begin{itemize}
    \item \textbf{Design Tokens:} Translucent cards with a 12px blur effect and soft linear gradients.
    \item \textbf{Utility Features:} Includes "Instant Predict" scorecards, "Intelligence EDA" for training data distribution, and a "Batch Processing Utility" for CSV files.
\end{itemize}

\section{System Resource Benchmarks}
\begin{table}[h]
\centering
\begin{tabular}{p{5cm}p{8cm}}
\toprule
\textbf{Metric} & \textbf{Measurement} \\
\midrule
Peak RAM Consumption & $\sim$6.2 GB (during BART operations). \\
Raw CPU Latency & 16,727ms (without pruning). \\
\textbf{Optimized Latency} & \textbf{1,095ms -- 1,400ms}. \\
\bottomrule
\end{tabular}
\caption{Computational Resource Benchmarks}
\end{table}

\chapter{Conclusion \& Future Work}
The Agentic Sentiment Intelligence System proves that context-aware agents can reliably solve complex sentiment problems like sarcasm. The project successfully achieved >87\% accuracy on complex datasets, proving agentic frameworks are superior to "black-box" classifiers.

\section{Future Roadmap}
\begin{itemize}
    \item \textbf{Multimodal Sentiment:} Incorporating image analysis (broken product photos).
    \item \textbf{Aspect-Based Analysis (ABSA):} Identifying specific features (e.g., Battery Life vs Screen Quality).
    \item \textbf{Multilingual Expansion:} Fine-tuning for regional Indian shorthand and mixed languages.
\end{itemize}

\appendix
\chapter{Interpretable AI \& Technical Q\&A}

\section{Sentence-Level Explainability via Prototype Reasoning}
Prototype-based reasoning provides sentence-level explainability by moving beyond the limitations of traditional word-level "black-box" interpretations, such as SHAP or LIME, which often struggle to explain complex intent like sarcasm when clear sentiment words are absent. This framework achieves intrinsic interpretability through the following mechanisms:

\begin{enumerate}
    \item \textbf{Learning Latent Clusters (Prototypes):} The model uses a specialized training regimen to learn a collection of prototype tensors, which encapsulate latent clusters of actual training samples. These prototypes represent diverse aspects of a class, such as specific ways sarcasm is expressed through analogy or government-skeptical remarks.
    \item \textbf{Similarity-Based Inference:} At the point of inference, the model does not rely on abstract weights but makes classification judgments based solely on the similarity between the input text and these learned prototypes. This similarity is often calculated using a Radial Basis Function (RBF) kernel, which measures the distance between the input's semantic embedding and the prototype vectors.
    \item \textbf{Explanation by Reference:} Because the final prediction is directly attributed to individual prototypes, the system can explain its reasoning by referencing the most closely aligned training examples. This process, known as prototype projection, replaces the abstract mathematical tensor with the nearest actual sentence from the training set, making the interpretation intuitive for humans.
    \item \textbf{Sentence-Level Context vs. Word-Level Attribution:} Unlike SHAP values, which assign individual weights to tokens and may result in non-intuitive score distributions across a sentence, prototype-based reasoning provides a holistic, human-readable comparison. For example, in sarcasm detection, a prototype model might explain a prediction by showing that a new input (e.g., about government regulations) is semantically similar to a known sarcastic training example about a similar topic.
\end{enumerate}

This "white-box" approach ensures that the model's decision process is transparent and verifiable by highlighting the specific training cases that most influenced its "thinking" at a sentence level.

\section{Handling Ironic Dissonance during Inference}
The model handles "ironic dissonance"---defined as a conflict between a low star rating and positive intensifiers like "Wow" or "Impressive"---through a specialized multi-tiered inference pipeline designed to understand context rather than just individual keywords.

The system processes these complex reviews in three distinct stages:
\begin{enumerate}
    \item \textbf{Linguistic Normalization (LQN):} Before the sentiment is analyzed, the Linguistic Quality Node (LQN) acts as a neural normalizer. It uses a transformer model (\texttt{pszemraj/grammar-synthesis-base}) to rewrite messy or slang-heavy reviews into standardized English. This ensures that the downstream specialist engine receives a clean sequence of words, making it easier to identify the structural markers of irony.
    \item \textbf{Context-Aware Specialist Analysis:} The core classification is handled by a fine-tuned DistilRoBERTa Specialist engine. This model was specifically trained on a dataset created through Active Sarcasm Mining, where 336 "Pure Sarcasm" cases (identified by their ironic dissonance) were given a 10x weighting to teach the model these rare linguistic patterns.
    \item \textbf{Symbolic Reasoning and Explanation (RE):} After the neural prediction is made, the Reasoning Engine (RE) provides a layer of Explainable AI (XAI). If the model identifies a review with ironic dissonance, the RE performs conflict detection and human-readable justification. It maps the neural confidence scores to business insights, explaining that the model prioritized specific modifiers over individual positive words to reach its negative conclusion.
\end{enumerate}

\section{Fixing E-commerce Slang and Typos via Neural Normalization}
Neural normalization fixes e-commerce slang and typos by acting as an "intelligence pipeline" that transforms messy, phonetic, or unstructured user input into standardized English before it reaches a sentiment classification engine. 

The LQN specifically addresses these challenges as follows:
\begin{enumerate}
    \item \textbf{Neural Rewriting of Phonetic Slang:} Instead of relying on simple static mappings, the LQN uses neural reasoning to understand the intended meaning of phonetic shorthand. For example, it can programmatically rewrite the input "i absolutli lik dis prduct" into "I absolutely like this product".
    \item \textbf{Overcoming "Neural Echoing" and Logic Failures:} Research conducted during this project revealed that simpler models suffered from "Neural Echoing"---merely repeating the input without fixing mid-word errors. To solve this, the final architecture implements Repetition Penalties (1.5 - 2.0) and Beam Search (k=5) to ensure high-accuracy denoising even in CPU-bound environments.
    \item \textbf{Intricate Data Refinery:} Beyond the neural model, the process includes specialized refinery steps like Emoji Demojization (converting emojis to \texttt{:smiling\_face:}) and Instruction Scaffolding to align model-specific task weights.
\end{enumerate}

\section{Prototype-based Reasoning for Sarcastic Analogies}
Sarcastic analogies often lack strong sentiment cues, which causes word-level explainers like SHAP to assign vague, distributed importance scores. Prototype-based reasoning addresses this through:
\begin{itemize}
    \item \textbf{Learning Latent Clusters:} Allocating a fixed number of prototypes initialized using k-means clustering on the training data.
    \item \textbf{Explanation by Reference (Prototype Projection):} Visualizing its reasoning by projecting a mathematical prototype vector onto its closest real-world datapoint in the training set.
    \item \textbf{Handling Analogies:} The model successfully matches inputs to the thematically closest training examples (e.g., matching a new analogy about government regulations to other known sarcastic comments about lawmakers).
    \item \textbf{Incongruity Loss:} Directly capturing the disparity between literal praise and underlying critical context.
\end{itemize}

\section{Identifying Sentiment Disparities via Incongruity Loss}
Incongruity loss identifies sentiment disparities by explicitly capturing the contradiction between the literal (explicit) and intended (implicit) meanings of a sentence.
\begin{enumerate}
    \item \textbf{Sentiment Encoding:} Decomposing text into the Explicit Part (sentiment words like "Great") and the Implicit Part (context clues like "phone broke in one day").
    \item \textbf{Capturing the Disparity:} Using cross-entropy to measure the difference between predicted sentiments of these parts.
    \item \textbf{Encouraging Disagreement:} For sarcastic labels, the loss specifically encourages disagreement. It learns that positive "great" + negative "broke" = Negative final label.
\end{enumerate}

\section{Integration with DistilRoBERTa Specialist}
Integrating incongruity loss with DistilRoBERTa involves transforming the training objective to a dual-stream analysis:
\begin{itemize}
    \item \textbf{Dual-Stream Encoding:} Feeding both explicit and implicit segments into the specialist separately.
    \item \textbf{Loss Function Penalty:} Mathematically penalizing the model when it fails to recognize the contradiction during fine-tuning.
    \item \textbf{Synergy with Active Data Mining:} Utilizing the Ironic Dissonance Miner to provide the ideal training samples for the incongruity loss to optimize against.
\end{itemize}

\section{Explicit vs. Implicit Sentiment in Sarcasm}
Explicit sentiment refers to surface-level words like "Wow" or "Impressive." Implicit sentiment is the underlying meaning derived from context (e.g., "...broke in one day"). Sarcasm occurs when the Explicit Part is positive but the Implicit Part is negative. Our "Neural Specialist" (DistilRoBERTa) maintains Contextual Persistence to prioritize the implicit context over surface-level praise.

\backmatter
\chapter{References}
\begin{enumerate}[label={[\arabic*]}]
\item Maragheh, R. Y., et al. (2025). ARAG: Agentic Retrieval Augmented Generation for Personalized Recommendation. \textit{arXiv:2506.21931}.
\item Cao, S., \& Hu, M. (2026). A Solicit-Then-Suggest Model of Agentic Purchasing. \textit{arXiv:2603.20972}.
\item Ayanami0730. (2026). A-RAG: Scaling Agentic Retrieval-Augmented Generation via Hierarchical Retrieval Interfaces. \textit{GitHub Repository}.
\item Scott, T. L., Goh, W. W., \& Khan, N. A. (2025). Aspect-Based Sentiment Analysis on Amazon Product Reviews Using a Novel Hybrid Machine Learning Algorithm. \textit{Journal of Universal Computer Science}.
\item Wen, X., \& Rezapour, R. (2026). A Transformer and Prototype-based Interpretable Model for Contextual Sarcasm Detection. \textit{ACL Anthology / WASSA 2026}.
\item Daraghmi, E., \& Zyadeh, N. (2026). Sentiment Classification of Amazon Product Reviews Based on Machine and Deep Learning Techniques: A Comparative Study. \textit{Future Internet}.
\item Thogesan, T., Nugaliyadde, A., \& Wong, K. W. (2025). Integration of Explainable AI Techniques with Large Language Models for Enhanced Interpretability for Sentiment Analysis. \textit{arXiv:2503.11948}.
\item Dubey, P., Dubey, P., \& Bokoro, P. N. (2025). Unpacking Sarcasm: A Contextual and Transformer-Based Approach for Improved Detection. \textit{Computers}.
\item Akter, P., et al. (2025). Sentiment Analysis of Consumer Feedback and Its Impact on Business Strategies by Machine Learning. \textit{The American Journal of Applied Sciences}.
\item Agrawal, K. (2025). The Future of E-Commerce with Agentic AI: Personalized, Task-Oriented AI for a Smarter Shopping Experience. \textit{OARJ Engineering and Technology}.
\item Trivedi, B. (2025). Explainable Agentic AI: Transforming E-Commerce Search With Transparency. \textit{Journal of International Crisis and Risk Communication Research}.
\item Chauhan, S., \& Kumar, A. (2024). A Transformer-based Model for Detecting Multilingual Sarcasm in Social Media Posts. \textit{CEUR Workshop Proceedings}.
\item Kaseb, A., \& Farouk, M. (2023). SAIDS: A Novel Approach for Sentiment Analysis Informed of Dialect and Sarcasm. \textit{arXiv:2301.02521}.
\item Abu Farha, I., et al. (2022). Sarcasm Detection is Way Too Easy! An Empirical Comparison of Human and Machine Sarcasm Detection. \textit{ACL Anthology / EMNLP Findings}.
\item Romio, G. (2026). 15 Best AI Sentiment Analysis Tools \& Use Cases in 2026. \textit{CloudTalk Blog}.
\end{enumerate}

\end{document}